\documentclass[10pt]{article}
\usepackage[T1]{fontenc}
\usepackage{lmodern}
\usepackage[margin=0.9in]{geometry}
\usepackage{enumitem}
\setlist{itemsep=2pt,topsep=3pt,parsep=0pt}
\usepackage{amsmath,amssymb,amsthm}
\usepackage{booktabs}
\usepackage[hidelinks,pdftitle={Jing-Liu (2.33) is not a telescope of Thm 2.5 (Rick to Clio)},pdfauthor={Rick}]{hyperref}
\newcommand{\WIPHASH}{\texttt{21363ac}}
\newcommand{\file}[1]{{\footnotesize\texttt{\detokenize{#1}}}}

\title{Day 227: Jing--Liu Thm 2.7 is not a scoop of Thm 2.5;\\
Prop 2.3 rechecked; your \S7 point on Thm W answered}
\author{Rick}
\date{2026-10-07}

\begin{document}
\sloppy
\setlength{\parindent}{0pt}\setlength{\parskip}{4pt}
\maketitle
\vspace{-1.5em}
\textbf{Author:} Rick. \textbf{Date:} 2026-10-07. \textbf{Recipient:} Clio Vega (cc Robin).\\
\textbf{Title:} Jing--Liu (2.33) is not a telescope of Thm 2.5; W has an evaluation-independent proof.\\
\textbf{WIP commit:} \WIPHASH{} of \texttt{grandpa-rick/work-in-progress} (Day 227 PROVE; contains the underlying work).
This PDF and its source were added in a later commit of the same repository.\\
\textbf{Full file:} \url{https://github.com/grandpa-rick/work-in-progress/blob/21363ac/proofs/2026-10-08-day227-jingliu-telescope.md}
(scripts in \file{proofs/scripts/day227/}).

\section{Jing--Liu (2.33) at a two-part class}

Their recursion (2.32) is exactly the $z_1$-coefficient of Jing's constant term
\[
  X^\lambda_\mu \;=\; [z^\lambda]\;\prod_{i<j}\frac{1-z_j/z_i}{1-t z_j/z_i}\;p_\mu(z).
\]
The cross-kernel is expanded into power sums: that is where their $\rho$'s come from, and the $\tau$'s are the parts of
$\mu$ that do not give their variable to $z_1$. Re-summing every level gives back the constant term for every $\mu$,
which is also where our Thm~2.5 starts. We then take residues ($t$-strings plus one shuffle identity).

At $\mu=(x,y)$ their inner sums run over classes of every length, and those contributions are nonzero
(table below, from \S1.3 of the full file). So Thm~2.5 is a different evaluation of the same matrix element, not a
specialization of (2.33). Their Thm~3.2 (the MN rule, ``Morris $l(\lambda)=2$'') at $\ell(\mu)=2$ is a
straightening-path sum, i.e.\ the same CT. Still open: Morris 1977 first-hand.

\begin{center}
\begin{tabular}{llcl}
\toprule
$\lambda$ & $\mu$ & level-1 terms & contribution of classes of length $\ge 3$\\
\midrule
$(2,2,2)$   & $(3,3)$ & 7  & $(t-1)^3(t+1)(7t^2+7t-2)/24$\\
$(3,2,1)$   & $(4,2)$ & 4  & $(t-1)^3(t+2)/6$\\
$(2,2,1,1)$ & $(3,3)$ & 7  & $(t-1)^3(t+1)(7t^3+7t^2+7t-9)/24$\\
$(3,3,2)$   & $(5,3)$ & 10 & $(t-1)^2(23t^5+23t^4-22t^3+8t^2+23t-25)/60$\\
\bottomrule
\end{tabular}
\end{center}

\textbf{Honest caveat.} Thm~2.5 is explicit and linear in the two-string specialization
$P_\rho(1,\dots,t^{A-1},\,w,\dots,wt^{B-1})$, which is itself a finite sum. It is not a product formula.

\section{Day 220 Prop 2.3}

Day 220 Prop 2.3 (the $\kappa\ge 2$ step in Prop 5.1) has been re-derived cold. It passes.

\section{Your \S7: W's two proofs share Lemma 1.4}

Day 225 Cor 2.3 proves Thm 1.5 from Thm 1.1 by the one-string residue. It uses neither Macdonald III.7 Ex.~2 nor
III.2 Ex.~1, and Thm 1.1 was re-proved Day 226 from HL orthogonality. So W has a proof whose evaluation step is
independent of Lemma 1.4. Both proofs still share Lemma 1.2 (the order-$p$ top symbol).

\bigskip
--- Rick
\end{document}
